\subsubsection*{Solution}
The expression belongs to the finite-dimensional Petz quasi-entropy framework introduced by \href{\detokenize{https://math.bme.hu/~petz/pdf/32quasi.pdf}}{Petz}; integral representations of the associated operator-convex functions are discussed by \href{\detokenize{https://arxiv.org/abs/1008.2529}}{Hiai, Mosonyi, Petz, and Bény}.

\paragraph*{Step-by-Step Derivation}

Define the Hermitian perturbation

\[
\Delta:=\rho-\sigma.
\]

Then

\[
\rho_t=\sigma+t\Delta, \qquad \rho_t-\sigma=t\Delta.
\]

Since left multiplication is linear,

\[
L_{\rho_t}=L_\sigma+tL_\Delta.
\]

For each $s>0$, introduce the superoperator

\[
A_t(s):=L_{\rho_t}+sR_\sigma =L_\sigma+tL_\Delta+sR_\sigma.
\]

The divergence along the path is therefore

\[
D_f^{\mathrm{std}}(\rho_t\|\sigma) = \int_0^\infty t^2\, \operatorname{tr} \left[ \Delta A_t(s)^{-1}(\Delta) \right]d\mu(s).
\]

\subparagraph*{1. Derivative of the inverse superoperator}

Because

\[
\frac{dA_t(s)}{dt}=L_\Delta,
\]

the inverse-derivative identity gives

\[
\frac{d}{dt}A_t(s)^{-1} = -A_t(s)^{-1}L_\Delta A_t(s)^{-1}.
\]

This follows directly from differentiating

\[
A_t(s)^{-1}A_t(s)=I.
\]

Consequently,

\[
\adjustbox{max width=\linewidth}{$\displaystyle \begin{aligned}
\frac{d}{dt}
\operatorname{tr}\left[
\Delta A_t(s)^{-1}(\Delta)
\right]
&=
\operatorname{tr}\left[
\Delta\frac{dA_t(s)^{-1}}{dt}(\Delta)
\right]
\\
&=
-\operatorname{tr}\left[
\Delta A_t(s)^{-1}
L_\Delta
A_t(s)^{-1}(\Delta)
\right].
\end{aligned}$}
\]

\subparagraph*{2. Differentiate the full integrand}

Applying the product rule,

\[
\adjustbox{max width=\linewidth}{$\displaystyle \begin{aligned}
\frac{d}{dt}
D_f^{\mathrm{std}}(\rho_t\|\sigma)
=
\int_0^\infty
\Bigl\{
&2t\,
\operatorname{tr}\left[
\Delta A_t(s)^{-1}(\Delta)
\right]
\\
&-t^2\,
\operatorname{tr}\left[
\Delta A_t(s)^{-1}
L_\Delta
A_t(s)^{-1}(\Delta)
\right]
\Bigr\}d\mu(s).
\end{aligned}$}
\]

\subparagraph*{3. Evaluate at $t=0.5$}

At $t=\frac12$,

\[
\rho_{1/2}=\frac{\rho+\sigma}{2}.
\]

Define

\[
B_s := L_{(\rho+\sigma)/2}+sR_\sigma.
\]

Substituting $2t=1$ and $t^2=\frac14$ gives

\[
\adjustbox{max width=\linewidth}{$\displaystyle \begin{aligned}
\left.
\frac{d}{dt}D_f^{\mathrm{std}}(\rho_t\|\sigma)
\right|_{t=1/2}
=
\int_0^\infty
\Bigg[
&
\operatorname{tr}\left[
\Delta B_s^{-1}(\Delta)
\right]
\\
&-\frac14
\operatorname{tr}\left[
\Delta B_s^{-1}
L_\Delta
B_s^{-1}(\Delta)
\right]
\Bigg]d\mu(s).
\end{aligned}$}
\]

\subparagraph*{4. Equivalent resolvent form}

Let

\[
A_{0,s}:=L_\sigma+sR_\sigma.
\]

Because

\[
B_s=A_{0,s}+\frac12L_\Delta,
\]

we have

\[
L_\Delta=2(B_s-A_{0,s}).
\]

Thus,

\[
B_s^{-1}L_\Delta B_s^{-1} = 2\left( B_s^{-1}-B_s^{-1}A_{0,s}B_s^{-1} \right).
\]

Substitution yields the equivalent expression

\[
{ \left. \frac{d}{dt}D_f^{\mathrm{std}}(\rho_t\|\sigma) \right|_{t=1/2} = \frac12\int_0^\infty \operatorname{tr}\left[ \Delta \left( B_s^{-1} + B_s^{-1}A_{0,s}B_s^{-1} \right)(\Delta) \right]d\mu(s) }.
\]

\subparagraph*{Differentiation under the integral}

In finite dimensions, a simple sufficient condition is that $\sigma$ be faithful:

\[
\sigma>0.
\]

In a neighborhood of $t=\frac12$,

\[
\rho_t=(1-t)\sigma+t\rho\geq(1-t)\sigma.
\]

Writing $m=\lambda_{\min}(\sigma)>0$, one obtains, for example for $\frac14\leq t\leq\frac34$,

\[
A_t(s)\geq m\left(\frac14+s\right)I,
\]

and hence

\[
\left\|A_t(s)^{-1}\right\| \leq \frac{4}{m(1+s)}.
\]

Both differentiated terms are consequently bounded by a constant times $(1+s)^{-1}$. The assumed condition

\[
\int_0^\infty\frac{1}{1+s}\,d\mu(s)<\infty
\]

then permits differentiation under the integral by dominated convergence. If $\sigma$ is singular, the same formula applies on the relevant support when the inverses are interpreted appropriately, but additional support conditions are needed.

No numerical scalar can be produced without specifying $\rho$, $\sigma$, and $\mu$.
